\textbf{Phần 1. Trắc nghiệm lựa chọn}
\setcounterpageref{ex}{0}

\begin{ex} %[5D4N21]
Kết quả của phép tính $7{,}8 \times 6$ là
\choice
{$42{,}8$}
{$46{,}4$}
{$48{,}8$}
{\True $46{,}8$}
\loigiai{
Ta có: $7{,}8 \times 6 = 46{,}8$.\\
\textbf{Đáp án của câu hỏi là D}
}
\end{ex}

\begin{ex} %[5D4N21]
Kết quả của phép tính $0{,}72 \times 50$ là
\choice
{\True $36$}
{$3{,}6$}
{$360$}
{$0{,}36$}
\loigiai{
Ta có: $0{,}72 \times 50 = 36$.\\
\textbf{Đáp án của câu hỏi là A}
}
\end{ex}

\begin{ex} %[5D4H21]
Kết quả của phép tính $4{,}3 \times 3{,}6$ là
\choice
{$154{,}8$}
{$1{,}548$}
{\True $15{,}48$}
{$15{,}44$}
\loigiai{
Hai thừa số có tổng cộng $2$ chữ số ở phần thập phân, vậy tích là $15{,}48$.\\
\textbf{Đáp án của câu hỏi là C}
}
\end{ex}

\begin{ex} %[5D4H21]
Kết quả của phép tính $7{,}5 \times 3{,}4$ là
\choice
{$24{,}5$}
{\True $25{,}5$}
{$25{,}50$}
{$255$}
\loigiai{
Ta có: $7{,}5 \times 3{,}4 = 25{,}5$.\\
\textbf{Đáp án của câu hỏi là B}
}
\end{ex}

\begin{ex} %[5D4H21]
Kết quả của phép tính $6{,}8 \times 0{,}52$ là
\choice
{\True $3{,}536$}
{$35{,}36$}
{$0{,}3536$}
{$353{,}6$}
\loigiai{
Hai thừa số có tổng cộng $3$ chữ số ở phần thập phân, vậy tích là $3{,}536$.\\
\textbf{Đáp án của câu hỏi là A}
}
\end{ex}

\begin{ex} %[5D4H21]
Cho biết $64 \times 57 = 3\,648$. Không thực hiện phép tính, giá trị của $6{,}4 \times 0{,}57$ là
\choice
{$36{,}48$}
{\True $3{,}648$}
{$0{,}3648$}
{$364{,}8$}
\loigiai{
Thừa số $6{,}4$ có $1$ chữ số thập phân, $0{,}57$ có $2$ chữ số thập phân $\Rightarrow$ Tích có $1 + 2 = 3$ chữ số thập phân. Vậy $6{,}4 \times 0{,}57 = 3{,}648$.\\
\textbf{Đáp án của câu hỏi là B}
}
\end{ex}

\begin{ex} %[5D4H21]
Tính thuận tiện biểu thức $2{,}5 \times 8{,}6 \times 4$ ta được kết quả là
\choice
{$8{,}6$}
{$860$}
{$86{,}4$}
{\True $86$}
\loigiai{
Ta có: $(2{,}5 \times 4) \times 8{,}6 = 10 \times 8{,}6 = 86$.\\
\textbf{Đáp án của câu hỏi là D}
}
\end{ex}

\begin{ex} %[5D4H21]
Tính thuận tiện biểu thức $6{,}84 \times 0{,}2 \times 5$ ta được kết quả là
\choice
{$68{,}4$}
{$0{,}684$}
{\True $6{,}84$}
{$13{,}68$}
\loigiai{
Ta có: $6{,}84 \times (0{,}2 \times 5) = 6{,}84 \times 1 = 6{,}84$.\\
\textbf{Đáp án của câu hỏi là C}
}
\end{ex}

\begin{ex} %[5D4V21]
Một ô tô chạy trên đường cao tốc, mỗi giờ đi được $84{,}5\text{ km}$. Trong $1{,}2$ giờ ô tô đó đi được quãng đường là
\choice
{$91{,}4\text{ km}$}
{\True $101{,}4\text{ km}$}
{$100{,}4\text{ km}$}
{$102{,}4\text{ km}$}
\loigiai{
Quãng đường ô tô đi được là: $84{,}5 \times 1{,}2 = 101{,}4\text{ (km)}$.\\
\textbf{Đáp án của câu hỏi là B}
}
\end{ex}

\begin{ex} %[5D4V21]
Biết rằng xay xát $1\text{ kg}$ thóc thì thu được $0{,}64\text{ kg}$ gạo. Hỏi xay xát $50\text{ kg}$ thóc loại đó thì thu được bao nhiêu ki-lô-gam gạo?
\choice
{$3{,}2\text{ kg}$}
{$30\text{ kg}$}
{\True $32\text{ kg}$}
{$320\text{ kg}$}
\loigiai{
Lượng gạo thu được là: $50 \times 0{,}64 = 32\text{ (kg)}$.\\
\textbf{Đáp án của câu hỏi là C}
}
\end{ex}

\begin{ex} %[5D4V21]
Mỗi cốc nước cam chứa $0{,}25\text{ l}$. Ba bạn nhỏ mỗi bạn uống một cốc nước cam thì uống hết tất cả số lít nước cam là
\choice
{$0{,}5\text{ l}$}
{$1\text{ l}$}
{$0{,}65\text{ l}$}
{\True $0{,}75\text{ l}$}
\loigiai{
Số lít nước cam cả ba bạn uống là: $0{,}25 \times 3 = 0{,}75\text{ (l)}$.\\
\textbf{Đáp án của câu hỏi là D}
}
\end{ex}

\begin{ex} %[5D4V21]
Mẹ của Mai mua $3\text{ kg}$ xoài hết $75\,000$ đồng. Cô Hà mua $3{,}5\text{ kg}$ xoài cùng loại và đưa cho người bán hàng tờ tiền $100\,000$ đồng. Người bán hàng phải trả lại cô Hà số tiền là
\choice
{\True $12\,500$ đồng}
{$15\,000$ đồng}
{$17\,500$ đồng}
{$87\,500$ đồng}
\loigiai{
Giá tiền $1\text{ kg}$ xoài là: $75\,000 : 3 = 25\,000$ (đồng).\\
Số tiền cô Hà mua $3{,}5\text{ kg}$ xoài là: $25\,000 \times 3{,}5 = 87\,500$ (đồng).\\
Số tiền trả lại là: $100\,000 - 87\,500 = 12\,500$ (đồng).\\
\textbf{Đáp án của câu hỏi là A}
}
\end{ex}

\textbf{Phần 2. Trắc nghiệm đúng sai}
\setcounterpageref{ex}{0}

\begin{ex} %[5D4H21]
Cho biết $3{,}6 \times 2{,}4 = 8{,}64$. Xét tính đúng sai của các khẳng định sau:
\choiceTF
{\True Phép tính $3{,}6 \times 24 = 86{,}4$.}
{Phép tính $36 \times 0{,}24 = 86{,}4$.}
{\True Phép tính $0{,}36 \times 2{,}4 = 0{,}864$.}
{\True Biểu thức $(5{,}3 \times 1{,}5) \times 2 = 5{,}3 \times (1{,}5 \times 2) = 15{,}9$.}
\loigiai{
\begin{itemize}
    \item $3{,}6 \times 24 = 86{,}4$ (thừa số thứ hai tăng 10 lần $\Rightarrow$ tích tăng 10 lần).
    \item $36 \times 0{,}24 = 8{,}64 \neq 86{,}4$.
    \item $0{,}36 \times 2{,}4 = 0{,}864$ (tích có 3 chữ số thập phân).
    \item $(5{,}3 \times 1{,}5) \times 2 = 5{,}3 \times 3 = 15{,}9$.
\end{itemize}
\begin{itemize}
    \item a) \textbf{Đúng}.
    \item b) \textbf{Sai}: Kết quả là $8{,}64$.
    \item c) \textbf{Đúng}.
    \item d) \textbf{Đúng}.
\end{itemize}
}
\end{ex}

\begin{ex} %[5D4V21]
Một căn phòng hình chữ nhật có chiều dài $4{,}3\text{ m}$ và chiều rộng $3{,}6\text{ m}$. Bác thợ dự định dùng loại sơn phủ sàn có giá $120\,000$ đồng/$1\text{ m}^2$. Xét tính đúng sai của các khẳng định sau:
\choiceTF
{\True Diện tích của căn phòng là $15{,}48\text{ m}^2$.}
{Chu vi của căn phòng là $7{,}9\text{ m}$.}
{\True Làm tròn diện tích căn phòng đến hàng đơn vị ta được $15\text{ m}^2$.}
{\True Chi phí tiền sơn sàn phòng là $1\,857\,600$ đồng.}
\loigiai{
\begin{itemize}
    \item Diện tích phòng: $4{,}3 \times 3{,}6 = 15{,}48\text{ m}^2$.
    \item Chu vi phòng: $(4{,}3 + 3{,}6) \times 2 = 7{,}9 \times 2 = 15{,}8\text{ m} \neq 7{,}9\text{ m}$.
    \item Làm tròn $15{,}48$ đến hàng đơn vị (hàng phần mười là $4 < 5$) được $15\text{ m}^2$.
    \item Tiền sơn: $15{,}48 \times 120\,000 = 1\,857\,600$ đồng.
\end{itemize}
\begin{itemize}
    \item a) \textbf{Đúng}.
    \item b) \textbf{Sai}: Chu vi là $15{,}8\text{ m}$.
    \item c) \textbf{Đúng}.
    \item d) \textbf{Đúng}.
\end{itemize}
}
\end{ex}

\textbf{Phần 3. Trả lời ngắn}
\setcounterpageref{ex}{0}

\begin{ex} %[5D4H21]
Tính kết quả phép nhân: $5{,}4 \times 39$. (Nhập kết quả dưới dạng số thập phân).
\shortans{210,6}
\loigiai{
\textbf{Đáp số:} $210{,}6$\\
$5{,}4 \times 39 = 210{,}6$.
}
\end{ex}

\begin{ex} %[5D4H21]
Tính kết quả phép nhân: $21{,}9 \times 5{,}1$. (Nhập kết quả dưới dạng số thập phân).
\shortans{111,69}
\loigiai{
\textbf{Đáp số:} $111{,}69$\\
$21{,}9 \times 5{,}1 = 111{,}69$.
}
\end{ex}

\begin{ex} %[5D4H21]
Tính giá trị biểu thức: $2{,}5 \times 8{,}6 \times 4$.
\shortans{86}
\loigiai{
\textbf{Đáp số:} $86$\\
$(2{,}5 \times 4) \times 8{,}6 = 10 \times 8{,}6 = 86$.
}
\end{ex}

\begin{ex} %[5D4V21]
Một ô tô đi quãng đường với vận tốc $84{,}5\text{ km/h}$. Hỏi sau $1{,}2$ giờ ô tô đi được bao nhiêu ki-lô-mét? (Nhập kết quả dưới dạng số thập phân).
\shortans{101,4}
\loigiai{
\textbf{Đáp số:} $101{,}4$\\
$84{,}5 \times 1{,}2 = 101{,}4\text{ (km)}$.
}
\end{ex}

\begin{ex} %[5D4V21]
Biết $1\text{ kg}$ thóc xay xát được $0{,}64\text{ kg}$ gạo. Hỏi $50\text{ kg}$ thóc xay xát được bao nhiêu ki-lô-gam gạo?
\shortans{32}
\loigiai{
\textbf{Đáp số:} $32$\\
$50 \times 0{,}64 = 32\text{ (kg)}$.
}
\end{ex}

\begin{ex} %[5D4V21]
Cô Hà mua $3{,}5\text{ kg}$ xoài giá $25\,000$ đồng/$1\text{ kg}$ và đưa tờ tiền $100\,000$ đồng. Hỏi người bán hàng phải trả lại bao nhiêu nghìn đồng? (Nhập kết quả dưới dạng số thập phân).
\shortans{12,5}
\loigiai{
\textbf{Đáp số:} $12{,}5$\\
Số tiền mua: $3{,}5 \times 25\,000 = 87\,500$ đồng. Số tiền trả lại: $100\,000 - 87\,500 = 12\,500$ đồng $= 12{,}5$ nghìn đồng.
}
\end{ex}
